\section{LLM-Based World Knowledge for Physical AI}
\label{sec:world-knowledge}

\emph{LLM-based world knowledge} refers to language-mediated world priors stored as parametric regularities in large language models~\citep{petroni2019language,roberts2020knowledge}.
Prompting and context expose parts of these priors as task-level hypotheses about objects, goals, actions, procedures, and constraints for Physical AI~\citep{brown2020language,wei2022chain}.

\subsection{World Knowledge as Coupled Priors}

World knowledge is better treated as a set of coupled priors than as a monolithic store: semantic knowledge grounds object categories and references; commonsense knowledge supplies default assumptions and safety constraints; procedural knowledge decomposes tasks; causal knowledge suggests outcomes and risks; spatial knowledge supports layout and manipulation preconditions; and affordance knowledge links objects or parts to possible interactions~\citep{arkin1990integrating,hagoort2004integration,schwenk2022okvqa}. 
In grounded Physical AI systems, spatial and affordance priors are often operationalized by models that connect linguistic relations to perceptual structure: recent spatial VLMs map relations such as ``inside'' or ``aligned with'' to metric or 3D structure~\citep{chen2024spatialvlm,yuan2024robopoint,song2025robospatial}. 
Affordance-oriented work similarly identifies which objects, parts, or regions support a requested action~\citep{qian2024affordancellm,huang2024manipvqa,chu2025threedaffordancellm}.

\subsection{Parametric World Knowledge in LLMs}

LLMs encode factual associations, procedural patterns, commonsense defaults, and task constraints through large-scale pretraining on corpora that repeatedly describe object functions, task procedures, safety constraints, and likely outcomes~\citep{petroni2019language,roberts2020knowledge,brown2020language,bommasani2021foundation}.
Because these descriptions recur across documents, they become parametric regularities whose strength tracks pretraining exposure~\citep{kandpal2023longtail,chang2024factualpretraining}.
Recent recall and materialization studies further probe which factual priors can be surfaced from LLMs~\citep{yuan2024factbench,hu2024gptkb}.
Embodied planning work provides a task-level diagnostic: LLMs can often retrieve plausible steps, object uses, and constraints for situated tasks~\citep{song2023llmplanner,rana2023sayplan,hua2024gensim2}.

\subsection{LLMs as Priors and Controllers in Physical AI}

LLMs contribute to Physical AI in four main ways.
For \textbf{task planning}, SayCan~\citep{ahn2022saycan} grounds LLM-generated step sequences with learned affordance values to filter physically infeasible actions, and Inner Monologue~\citep{huang2022innermonologue} enables closed-loop replanning by feeding back environmental observations as language~\citep{huang2022zeroshotplanners}.
Moving from prose to code, \textbf{skill generation} approaches such as Code as Policies~\citep{liang2023codeaspolicies} and ProgPrompt~\citep{singh2023progprompt} prompt LLMs to produce executable robot programs, while Voyager~\citep{wang2023voyager} uses LLMs to iteratively build open-ended skill libraries.
At the level of \textbf{goal and reward specification}, VoxPoser~\citep{huang2023voxposer} composes spatial value maps from language instructions, reducing manual reward engineering.
Finally, as \textbf{agentic orchestrators}, LLMs coordinate tools and sub-modules across multi-step tasks~\citep{rana2023sayplan,hua2024gensim2}, acting as closed-loop decision engines rather than one-shot planners.
Together, these uses establish LLMs as the semantic and procedural scaffolding of Physical AI pipelines---a foundation whose limits are examined next.

\subsection{Limits of Language-Only Physical Reasoning}

The systems reviewed above reveal a common pattern: LLMs are effective at proposing goals, plans, and constraints, but their outputs become more reliable when paired with grounding, verification, or action interfaces~\citep{ahn2022saycan,huang2022innermonologue,rana2023sayplan}. 
The core limitation is the abstraction gap between language and physics. 
Language compresses continuous physical states into sparse descriptions and often omits pose, velocity, contact, deformation, uncertainty, and embodiment constraints. 
In this sense, LLM-based priors remain linguistic: they may suggest that glass is fragile or that a handle can be grasped, but they do not estimate geometry, contact, force, friction, or future trajectories~\citep{qiu2025phybench,xu2025physense}. 
LLMs can therefore hallucinate objects, propose infeasible plans, or assume that actions succeed.
Recent planning analyses, including LLM-modulo studies and PlanBench updates, sharpen this point by showing that language-only models remain unreliable without external models or verifiers~\citep{kambhampati2024llmscantplan,valmeekam2024planbench}. 
Physical reasoning benchmarks further expose weaknesses in dynamics, vision-grounded physics, and tool use~\citep{xiang2025seephys,zhang2025phystoolbench}. 
These limitations motivate the remainder of this survey, which progressively addresses perceptual grounding, action grounding, and predictive world modeling.
LLMs are therefore useful interfaces and coordinators, but they form an important yet insufficient layer of Physical AI.
